\chapter{Краткое описание алгоритма FP-Growth}

В основе алгоритма FP-Growth (Frequent Pattern Growth) лежит построение FP-дерева (Frequent Pattern Tree), которое представляет транзакции базы данных в сжатом виде, сохраняя только частые элементы. Сначала алгоритм сканирует транзакции, определяет поддержку каждого элемента и отфильтровывает редкие элементы на основе минимального порога. Затем объекты в транзакциях сортируются по убыванию поддержки и добавляются в FP-дерево с объединением общих префиксов. Затем из дерева извлекаются наиболее часто встречающиеся наборы объектов. На основе полученных наборов формируются ассоциативные правила с учётом порога уверенности. Такой подход позволяет избежать многократного сканирования базы данных транзакций и генерации большого числа кандидатов, как в алгоритме Apriori.

\clearpage
